\usepackage[papersize={210mm,297mm},lmargin=2.5cm,rmargin=2.5cm,top=3.5cm,bottom=3.5cm]{geometry}

\renewcommand{\baselinestretch}{1.25}
\setlength{\parskip}{1mm}

\usepackage{scrextend}
\usepackage{comment}
\usepackage{caption}
\usepackage{subcaption}
\usepackage{amsmath,amsfonts,amssymb,amsthm,amscd}
\usepackage{pdfsync}
\usepackage{latexsym,xcolor,graphicx}
\usepackage[nottoc]{tocbibind}
\usepackage[utf8]{inputenc}					
\usepackage[T1]{fontenc}
\usepackage[spanish,es-lcroman,es-noquoting]{babel}
\usepackage{csquotes}
\usepackage{tikz}
\usepackage{tikz-cd}
\usepackage{appendix}
\usepackage[hidelinks]{hyperref}
\usepackage{enumitem}\setlist{nosep,after=\vspace{1mm},before=\vspace{0mm}}
\usepackage{graphicx,epstopdf}
\usepackage{tcolorbox}
\usepackage{amssymb,cancel}
\usepackage{amsfonts,amsmath,color}
\usepackage{amsthm,mathrsfs}
\usepackage{enumitem,multicol,bbding}
\usepackage[mathscr]{euscript}
\usepackage{fancyhdr}
\usepackage{etoolbox}
\usepackage{csquotes}
\usepackage{etoolbox}
\usepackage{scrextend}

\usepackage{etoolbox}
\newcommand{\zerodisplayskips}{%
  \setlength{\abovedisplayskip}{2mm}%
  \setlength{\belowdisplayskip}{2mm}%
  \setlength{\abovedisplayshortskip}{1mm}%
  \setlength{\belowdisplayshortskip}{1mm}}
\appto{\normalsize}{\zerodisplayskips}
\appto{\small}{\zerodisplayskips}
\appto{\footnotesize}{\zerodisplayskips}

\usepackage[style=ieee]{biblatex}
%\usepackage{biblatex}
\addbibresource{referencias.bib}

\usepackage{titlesec}
\titleformat{\chapter}[display]
{\bfseries\huge}
{\filleft\huge\chaptertitlename~\thechapter}
{0.5cm}
{\titlerule\filright}
[\titlerule]
\titlespacing*{\chapter}{0pt}{-1.5cm}{2cm}
\titlespacing*{\section}{0pt}{0.25cm}{0.25cm}

\patchcmd{\thebibliography}{\section*{\refname}}{}{}{}

% Cabeceira
\usepackage{fancyhdr}
\pagestyle{fancy}
\fancyhf{}
\renewcommand{\chaptermark}[1]{\markboth{\textbf{\thechapter. #1}}{}}% Formato para o capítulo: N. Nome
\renewcommand{\sectionmark}[1]{\markright{\textbf{\thesection. #1}}}% Formato para a sección: N.M. Nome
\fancyhead[LO]{\rightmark}% Nas páxinas impares, parte esquerda do encabezado, aparecerá o nome do capítulo
\fancyhead[RE]{\leftmark}% Nas páxinas pares, parte dereita do encabezado, aparecerá o nombre da sección
\fancyhead[RO,LE]{\textbf{\thepage}} % Números de páxina nas esquinas dos encabezados

\renewcommand{\quote}{\list{}{\rightmargin=\leftmargin\topsep=0pt}\item\relax}

\newcommand{\portada}[3]{
\thispagestyle{empty}
\definecolor{azulUSC}{RGB}{13,38,120}
\includegraphics[width=13cm]{logomath}
\vspace*{2cm}
\centerline{{\Large\bf  Trabajo Fin de Grado}}
\vspace{2.6cm}
{\center{\Huge\bfseries\color{azulUSC} #1}\par}		
\vspace{1cm}
{\Large
\centerline{#2}		
\vspace{6.4cm}
\centerline{\sf #3}        
\vspace{1.25cm}
\centerline{\sf UNIVERSIDADE DE SANTIAGO DE COMPOSTELA}
}
\enlargethispage{1cm}
\clearpage
}

\newcommand{\primeira}[3]{
\thispagestyle{empty}
\enlargethispage{1cm}
\begin{center}\Large
{\sf GRADO DE MATEM\'ATICAS}\\
\bigskip
{\bfseries Trabajo Fin de Grado}\\
\vspace{4cm}
{\bfseries\Huge #1}\\		
\vspace{1.2cm}
{#2}		\par		     
\end{center}
\vspace{6.5cm}
\begin{center}
\Large
{#3}	\par			    
\vspace{1cm}
{\sf UNIVERSIDADE DE SANTIAGO DE COMPOSTELA}
\end{center}}

\newtheorem{theorem}{Teorema}[chapter]
\newtheorem{proposition}[theorem]{Proposición}
\newtheorem{lemma}[theorem]{Lema}
\newtheorem{corollary}[theorem]{Corolario}
\newtheorem{algorithm}[theorem]{Algoritmo}
\newtheorem{axiom}[theorem]{Axioma}
\newtheorem{case}[theorem]{Caso}
\newtheorem{conclusion}[theorem]{Conclusión}
\newtheorem{condition}[theorem]{Condición}
\newtheorem{conjetura}[theorem]{Conjetura}				
\newtheorem{criterion}[theorem]{Criterio}
\newtheorem{ejercicio}[theorem]{Ejercicio}	
			
\theoremstyle{definition}
\newtheorem{definition}[theorem]{Definición}
\newtheorem{ejemplo}[theorem]{Ejemplo}				
\newtheorem{agradecimientos}[theorem]{Agradecimientos}	

\theoremstyle{remark}
\newtheorem{remark}[theorem]{Observación}
\newtheorem{notation}[theorem]{Notación}
\newtheorem{problem}[theorem]{Problema}
\newtheorem{solution}[theorem]{Solución}

\newtheorem*{remark*}{Observación}
\newtheorem*{remarkimp*}{Observación importante}


\addto\captionsspanish{
	\renewcommand{\contentsname}{\'Indice} 
	\renewcommand\appendixname{Anexo}
	\renewcommand\appendixpagename{Anexos}
}


%PAQUETES E COMANDOS EXTRA
\theoremstyle{definition}
\newtheoremstyle{exampstyle}
{2mm} % Space above
{-2mm} % Space below
{} % Body font
{} % Indent amount
{\bfseries} % Theorem head font
{.} % Punctuation after theorem head
{.5em} % Space after theorem head
{} % Theorem head spec (can be left empty, meaning `normal')
\theoremstyle{exampstyle}
\newtheorem{innercustomgeneric}{\customgenericname}
\providecommand{\customgenericname}{}
\newcommand{\newcustomtheorem}[2]{%
  \newenvironment{#1}[1]
  {%
   \renewcommand\customgenericname{#2}%
   \renewcommand\theinnercustomgeneric{##1}%
   \innercustomgeneric
  }
  {\endinnercustomgeneric}
}
\newcustomtheorem{cdefinition}{Definición}

\AfterEndEnvironment{cdefinition}{\noindent\ignorespaces}

\def\N{\ensuremath{\mathbb{N}}} %-- la N de los naturales
\def\R{\ensuremath{\mathbb{R}}} %-- La R de los reales.
\def\C{\ensuremath{\mathbb{C}}} %-- La C de los complejos.
\def\Q{\ensuremath{\mathbb{Q}}}
\def\K{\ensuremath{\mathbb{K}}}

\newcommand{\vertiii}[1]{{\left\vert\kern-0.25ex\left\vert\kern-0.25ex\left\vert #1 
    \right\vert\kern-0.25ex\right\vert\kern-0.25ex\right\vert}}


% Definimos el entorno para el PDF (LaTeX)
% Cambiamos las viñetas por un contador numérico con el prefijo "N"
\newenvironment{lista-n}%
  {\begin{itemize}[label=\textbf{(N\arabic*)}, ref=(N\arabic*), leftmargin=2.5em]\setcounter{enumi}{0}}%
  {\end{itemize}}